\chapter{Introduction}

These four laboratory sessions were a hands-on journey through the core principles of communication systems, running in parallel with our lectures to bridge theory and practice in LabVIEW. From the outset, we completed every bonus exercise and extended several experiments beyond what was asked, driven by a genuine curiosity to explore behaviours the lab handouts only hinted at.

Throughout, we maintained a consistently rigorous mathematical approach by grounding every result in first-principles derivations, from SNR analysis under AWGN for AM demodulators, to Bessel function interpretations of FM spectra, matched filter theory for BPSK/DPSK, and QAM constellation trade-offs. Every LabVIEW block diagram was treated as a direct realisation of the underlying theory, with experimental observations systematically verified against analytical predictions.

From the Central Limit Theorem and AM/FM systems to live USRP transmissions and error-correction coding, each lab gave us a deeper intuition for the engineering decisions behind real communication systems.


